\chapter{2020年全国I卷（理）}

\section{单选题}

\begin{problem}
若 \(z=1+\i\)，则 \(\left\lvert z^2-2z \right\rvert=\)（\quad）

\choices
    {\(0\)}
    {\(1\)}
    {\(\sqrt{2}\)}
    {\(2\)}

\end{problem}

\begin{answer}
D
\end{answer}

\begin{solution}
若 \(z=1+\i\)，则
\(z^2-2z=\left( 1+\i \right)^2-2\left( 1+\i \right)=2\i-2-2\i=-2\)，
所以
\(\left\lvert z^2-2z \right\rvert=\left\lvert -2 \right\rvert=2\).
故选 D.
\end{solution}

\begin{problem}
设集合 \(A=\{x\mid x^2-4\le 0\}\)，\(B=\{x\mid 2x+a\le 0\}\)，且 \(A\cap B=\{x\mid -2\le x\le 1\}\)，则 \(a=\)（\quad）

\choices
    {\(-4\)}
    {\(-2\)}
    {\(2\)}
    {\(4\)}

\end{problem}

\begin{answer}
B
\end{answer}

\begin{solution}
由题意可得 \(A=\{x\mid -2\le x\le 2\}\)，\(B=\left\{ x\mid x\le -\frac{a}{2} \right\}\).
又因为
\(A\cap B=\{x\mid -2\le x\le 1\}\)，
所以
\(-\frac{a}{2}=1\).
故
\(a=-2\).
故选 B.
\end{solution}

\begin{problem}
埃及胡夫金字塔是古代世界建筑奇迹之一，它的形状可视为一个正四棱锥. 以该四棱锥的高为边长的正方形面积等于该四棱锥一个侧面三角形的面积，则其侧面三角形底边上的高与底面正方形的边长的比值为（\quad）

\bitmapfigure[max width=6.6cm]{img/2020/national_paper_1_science/q3_fig1.png}

\choices
    {\(\dfrac{\sqrt{5}-1}{4}\)}
    {\(\dfrac{\sqrt{5}-1}{2}\)}
    {\(\dfrac{\sqrt{5}+1}{4}\)}
    {\(\dfrac{\sqrt{5}+1}{2}\)}

\end{problem}

\begin{answer}
C
\end{answer}

\begin{solution}
设正四棱锥的高为 \(h\)，底面边长为 \(a\)，侧面三角形底边上的高为 \(h'\).
由题意可得 \(h^2=\dfrac{1}{2}ah'\)，\(h^2=h'^2-\left( \dfrac{a}{2} \right)^2\).
所以
\(h'^2-\left( \dfrac{a}{2} \right)^2=\frac{1}{2}ah'\).
即
\(4\left( \frac{h'}{a} \right)^2-2\left( \frac{h'}{a} \right)-1=0\).
解得
\(\frac{h'}{a}=\frac{\sqrt{5}+1}{4}\)
（负值舍去）.
故选 C.
\end{solution}

\begin{problem}
已知 \(A\) 为抛物线 \(C:y^2=2px\)（\(p>0\)）上一点，点 \(A\) 到 \(C\) 的焦点的距离为 \(12\)，到 \(y\) 轴的距离为 \(9\)，则 \(p=\)（\quad）

\choices
    {\(2\)}
    {\(3\)}
    {\(6\)}
    {\(9\)}

\end{problem}

\begin{answer}
C
\end{answer}

\begin{solution}
因为抛物线上的点到焦点的距离等于到准线的距离，所以
\(9+\frac{p}{2}=12\).
解得
\(p=6\).
故选 C.
\end{solution}

\begin{problem}
某校一个课外学习小组为研究某作物种子的发芽率 \(y\) 和温度 \(x\)（单位：\(^\circ\)C）的关系，在 \(20\) 个不同的温度条件下进行种子发芽实验，由实验数据 \(\left( x_i,y_i \right)\)（\(i=1\)，\(2\)，\(\cdots\)，\(20\)）得到下面的散点图：

\bitmapfigure[max width=5.6cm]{img/2020/national_paper_1_science/q5_fig1.png}

由此散点图，在 \(10\)\(^\circ\)C 至 \(40\)\(^\circ\)C 之间，下面四个回归方程类型中最适宜作为发芽率 \(y\) 和温度 \(x\) 的回归方程类型的是（\quad）

\choices
    {\(y=a+bx\)}
    {\(y=a+bx^2\)}
    {\(y=a+b\e^x\)}
    {\(y=a+b\ln x\)}

\end{problem}

\begin{answer}
D
\end{answer}

\begin{solution}
由散点图可知，在 \(10\)\(^\circ\)C 至 \(40\)\(^\circ\)C 之间，发芽率 \(y\) 和温度 \(x\) 所对应的点 \(\left( x,y \right)\) 在一段对数函数的曲线附近.
结合选项可知，\(y=a+b\ln x\) 可作为发芽率 \(y\) 和温度 \(x\) 的回归方程类型.
故选 D.
\end{solution}

\begin{problem}
函数 \(f(x)=x^4-2x^3\) 的图象在点 \(\left( 1,f\left( 1 \right) \right)\) 处的切线方程为（\quad）

\choices
    {\(y=-2x-1\)}
    {\(y=-2x+1\)}
    {\(y=2x-3\)}
    {\(y=2x+1\)}

\end{problem}

\begin{answer}
B
\end{answer}

\begin{solution}
由
\(f(x)=x^4-2x^3\)
得
\(f'(x)=4x^3-6x^2\).
所以
\(f'(1)=4-6=-2\)，\(f(1)=1-2=-1\).
故切线方程为
\(y-\left( -1 \right)=-2\left( x-1 \right)\)，
即
\(y=-2x+1\).
故选 B.
\end{solution}

\begin{problem}
设函数 \(f(x)=\cos\left( \omega x+\dfrac{\pi}{6} \right)\) 在 \([-\pi,\pi]\) 的图象大致如图，则 \(f(x)\) 的最小正周期为（\quad）

\bitmapfigure[max width=0.7\linewidth]{img/2020/national_paper_1_science/q7_fig1.png}

\choices
    {\(\dfrac{10\pi}{9}\)}
    {\(\dfrac{7\pi}{6}\)}
    {\(\dfrac{4\pi}{3}\)}
    {\(\dfrac{3\pi}{2}\)}

\end{problem}

\begin{answer}
C
\end{answer}

\begin{solution}
由图象可得最小正周期小于
\(\pi-\left( -\frac{4\pi}{9} \right)=\frac{13\pi}{9}\)，
大于
\(2\left( \pi-\frac{4\pi}{9} \right)=\frac{10\pi}{9}\)，
故排除 \(A\)，\(D\).
又由图象可得
\(f\left( -\frac{4\pi}{9} \right)=\cos\left( -\frac{4\pi}{9}\omega+\frac{\pi}{6} \right)=0\)，
所以
\(-\frac{4\pi}{9}\omega+\frac{\pi}{6}=k\pi+\frac{\pi}{2}\)，\(k\in\Z\).
若选 \(B\)，则
\(\omega=\frac{2\pi}{\frac{7\pi}{6}}=\frac{12}{7}\)，
代入上式可得 \(k\) 不为整数，故排除 \(B\).
若选 \(C\)，则
\(\omega=\frac{2\pi}{\frac{4\pi}{3}}=\frac{3}{2}\)，
代入上式可得 \(k=-1\)，成立.
故选 C.
\end{solution}

\begin{problem}
\(\left( x+\dfrac{y^2}{x} \right)\left( x+y \right)^5\) 的展开式中 \(x^3y^3\) 的系数为（\quad）

\choices
    {\(5\)}
    {\(10\)}
    {\(15\)}
    {\(20\)}

\end{problem}

\begin{answer}
C
\end{answer}

\begin{solution}
因为
\(\left( x+\frac{y^2}{x} \right)\left( x+y \right)^5=\frac{\left( x^2+y^2 \right)\left( x+y \right)^5}{x}\)，
所以原展开式中 \(x^3y^3\) 的系数，就是
\(\left( x^2+y^2 \right)\left( x+y \right)^5\)
展开式中 \(x^4y^3\) 的系数.
展开式中含 \(x^4y^3\) 的项为
\(x^2\cdot C_5^2x^2y^3+y^2\cdot C_5^4x^4y=15x^4y^3\).
所以所求系数为 \(15\).
故选 C.
\end{solution}

\begin{problem}
已知 \(\alpha\in\left( 0,\pi \right)\)，且 \(3\cos 2\alpha-8\cos\alpha=5\)，则 \(\sin\alpha=\)（\quad）

\choices
    {\(\dfrac{\sqrt{5}}{3}\)}
    {\(\dfrac{2}{3}\)}
    {\(\dfrac{1}{3}\)}
    {\(\dfrac{\sqrt{5}}{9}\)}

\end{problem}

\begin{answer}
A
\end{answer}

\begin{solution}
由
\(3\cos 2\alpha-8\cos\alpha=5\)
得
\(3\left( 2\cos^2\alpha-1 \right)-8\cos\alpha-5=0\)，
即
\(3\cos^2\alpha-4\cos\alpha-4=0\).
解得
\(\cos\alpha=2\)
（舍去），或
\(\cos\alpha=-\frac{2}{3}\).
因为 \(\alpha\in\left( 0,\pi \right)\)，所以 \(\alpha\in\left( \frac{\pi}{2},\pi \right)\)，从而
\(\sin\alpha=\sqrt{1-\cos^2\alpha}=\sqrt{1-\left( -\frac{2}{3} \right)^2}=\frac{\sqrt{5}}{3}\).
故选 A.
\end{solution}

\begin{problem}
已知 \(A\)，\(B\)，\(C\) 为球 \(O\) 的球面上的三个点，\(\odot O_1\) 为 \(\triangle ABC\) 的外接圆. 若 \(\odot O_1\) 的面积为 \(4\pi\)，\(AB=BC=AC=OO_1\)，则球 \(O\) 的表面积为（\quad）

\choices
    {\(64\pi\)}
    {\(48\pi\)}
    {\(36\pi\)}
    {\(32\pi\)}

\end{problem}

\begin{answer}
A
\end{answer}

\begin{solution}
由题意可知图形如图：
\bitmapfigure[max width=5.2cm]{img/2020/national_paper_1_science/q10_solution_fig1.png}
因为 \(\odot O_1\) 的面积为 \(4\pi\)，所以
\(O_1A=2\).
又因为 \(\triangle ABC\) 为正三角形，所以
\(O_1A=AB\sin 60^\circ\).
故
\(AB=BC=AC=OO_1=2\sqrt{3}\).
设球 \(O\) 的半径为 \(R\)，则
\(R=\sqrt{AO_1^2+OO_1^2}=\sqrt{2^2+\left( 2\sqrt{3} \right)^2}=4\).
所以球 \(O\) 的表面积为
\(4\pi R^2=4\pi\times 4^2=64\pi\).
故选 A.
\end{solution}

\begin{problem}
已知 \(\odot M:x^2+y^2-2x-2y-2=0\)，直线 \(l:2x+y+2=0\)，\(P\) 为 \(l\) 上的动点. 过点 \(P\) 作 \(\odot M\) 的切线 \(PA\)，\(PB\)，切点为 \(A\)，\(B\)，当 \(\left\lvert PM \right\rvert\cdot \left\lvert AB \right\rvert\) 最小时，直线 \(AB\) 的方程为（\quad）

\choices
    {\(2x-y-1=0\)}
    {\(2x+y-1=0\)}
    {\(2x-y+1=0\)}
    {\(2x+y+1=0\)}

\end{problem}

\begin{answer}
D
\end{answer}

\begin{solution}
化圆 \(M\) 为
\(\left( x-1 \right)^2+\left( y-1 \right)^2=4\)，
所以圆心为 \(M\left( 1,1 \right)\)，半径为 \(2\).
由面积关系可得 \(S_{\text{四边形 }PAMB}=\frac{1}{2}\left\lvert PM \right\rvert\cdot \left\lvert AB \right\rvert=2S_{\triangle PAM}=\left\lvert PA \right\rvert\cdot \left\lvert AM \right\rvert=2\sqrt{\left\lvert PM \right\rvert^2-4}\).
因此要使 \(\left\lvert PM \right\rvert\cdot \left\lvert AB \right\rvert\) 最小，只需使 \(\left\lvert PM \right\rvert\) 最小，此时 \(PM\perp l\).
直线 \(PM\) 的方程为
\(y-1=\frac{1}{2}\left( x-1 \right)\)，
即
\(y=\frac{1}{2}x+\frac{1}{2}\).
联立
\(y=\dfrac{1}{2}x+\dfrac{1}{2}\) 和 \(2x+y+2=0\)，解得 \(P\left( -1,0 \right)\).
所以以 \(PM\) 为直径的圆的方程为
\(x^2+\left( y-\frac{1}{2} \right)^2=\frac{5}{4}\).
再与圆 \(M\) 联立，可得直线 \(AB\) 的方程为
\(2x+y+1=0\).
故选 D.
\end{solution}

\begin{problem}
若 \(2^a+\log_2 a=4^b+2\log_4 b\)，则（\quad）

\choices
    {\(a>2b\)}
    {\(a<2b\)}
    {\(a>b^2\)}
    {\(a<b^2\)}

\end{problem}

\begin{answer}
B
\end{answer}

\begin{solution}
因为
\(2^a+\log_2 a=4^b+2\log_4 b=2^{2b}+\log_2 b\).
又因为
\(2^{2b}+\log_2 b<2^{2b}+\log_2 2b\)，
所以
\(2^a+\log_2 a<2^{2b}+\log_2 2b\).
令
\(f(x)=2^x+\log_2 x\)，
则 \(f(x)\) 在 \(\left( 0,+\infty \right)\) 内单调递增.
因为
\(f(a)<f(2b)\)，
所以
\(a<2b\).
故选 B.
\end{solution}
\section{填空题}

\begin{problem}
若 \(x\)，\(y\) 满足约束条件 \(\begin{cases} 2x+y-2\le 0,\\ x-y-1\ge 0,\\ y+1\ge 0 \end{cases}\)，则 \(z=x+7y\) 的最大值为\fillinblank{}.

\end{problem}

\begin{answer}
1
\end{answer}

\begin{solution}
由 \(2x+y-2=0\) 和 \(x-y-1=0\) 联立，解得 \(A\left( 1,0 \right)\).
目标函数为
\(z=x+7y\)，
可化为
\(y=-\frac{1}{7}x+\frac{1}{7}z\).
当直线
\(y=-\frac{1}{7}x+\frac{1}{7}z\)
过点 \(A\) 时，在 \(y\) 轴上的截距最大，此时 \(z\) 取得最大值.
所以
\(z_{\max}=1+7\times 0=1\).
故答案为 \(1\).
\end{solution}

\begin{problem}
设 \(\bs a\)，\(\bs b\) 为单位向量，且 \(\left\lvert \bs a+\bs b \right\rvert=1\)，则 \(\left\lvert \bs a-\bs b \right\rvert=\fillinblank{}\).

\end{problem}

\begin{answer}
\(\sqrt{3}\)
\end{answer}

\begin{solution}
因为 \(\bs a\)，\(\bs b\) 为单位向量，且 \(\left\lvert \bs a+\bs b \right\rvert^2=1\)，所以 \(\left\lvert \bs a \right\rvert^2+2\bs a\cdot \bs b+\left\lvert \bs b \right\rvert^2=1\).
即
\(1+2\bs a\cdot \bs b+1=1\)，
故
\(2\bs a\cdot \bs b=-1\).
于是 \(\left\lvert \bs a-\bs b \right\rvert=\sqrt{\left\lvert \bs a \right\rvert^2-2\bs a\cdot \bs b+\left\lvert \bs b \right\rvert^2}=\sqrt{1-\left( -1 \right)+1}=\sqrt{3}\).
故答案为 \(\sqrt{3}\).
\end{solution}

\begin{problem}
已知 \(F\) 为双曲线 \(C:\dfrac{x^2}{a^2}-\dfrac{y^2}{b^2}=1\)（\(a>0\)，\(b>0\)）的右焦点，\(A\) 为 \(C\) 的右顶点，\(B\) 为 \(C\) 上的点，且 \(BF\perp x\) 轴. 若 \(AB\) 的斜率为 \(3\)，则 \(C\) 的离心率为\fillinblank{}.

\end{problem}

\begin{answer}
2
\end{answer}

\begin{solution}
设右焦点为 \(F\left( c,0 \right)\)，右顶点为 \(A\left( a,0 \right)\).
因为 \(B\) 在双曲线 \(C\) 上，且 \(BF\perp x\) 轴，所以
\(B\left( c,\frac{b^2}{a} \right)\).
由 \(AB\) 的斜率为 \(3\)，可得
\(\frac{\frac{b^2}{a}-0}{c-a}=3\).
又因为
\(b^2=c^2-a^2\)，
代入可得
\(c^2=3ac-2a^2\).
设双曲线的离心率为 \(e\)，则
\(e=\frac{c}{a}\)，
所以
\(e^2-3e+2=0\).
因为 \(e>1\)，所以
\(e=2\).
故答案为 \(2\).
\end{solution}

\begin{problem}
如图，在三棱锥 \(P-ABC\) 的平面展开图中，\(AC=1\)，\(AB=AD=\sqrt{3}\)，\(AB\perp AC\)，\(AB\perp AD\)，\(\angle CAE=30^\circ\)，则 \(\cos \angle FCB=\fillinblank{}\).

\bitmapfigure[max width=5.2cm]{img/2020/national_paper_1_science/q16_fig1.png}

\end{problem}

\begin{answer}
\(-\dfrac{1}{4}\)
\end{answer}

\begin{solution}
由已知可得
\(BD=\sqrt{2}AB=\sqrt{6}\)，\(BC=2\).
又因为 \(D\)，\(E\)，\(F\) 三点重合，所以
\(AE=AD=\sqrt{3}\)，\(BF=BD=\sqrt{6}\).
在 \(\triangle ACE\) 中，由余弦定理可得 \(CE^2=AC^2+AE^2-2AC\cdot AE\cos \angle CAE=1+3-2\times \sqrt{3}\times \frac{\sqrt{3}}{2}=1\).
故
\(CE=CF=1\).
于是 \(\cos \angle FCB=\frac{BC^2+CF^2-BF^2}{2BC\cdot CF}=\frac{4+1-6}{2\times 2\times 1}=-\frac{1}{4}\).
故答案为 \(-\dfrac{1}{4}\).
\end{solution}
\section{解答题}

\begin{problem}
设 \(\{a_n\}\) 是公比不为 \(1\) 的等比数列，\(a_1\) 为 \(a_2\)，\(a_3\) 的等差中项.

\begin{enumerate}
\item 求 \(\{a_n\}\) 的公比；

\item 若 \(a_1=1\)，求数列 \(\{na_n\}\) 的前 \(n\) 项和.
\end{enumerate}
\end{problem}

\begin{solution}
（1）设 \(\{a_n\}\) 是公比 \(q\) 不为 \(1\) 的等比数列.
由题意可得
\(2a_1=a_2+a_3\)，
即
\(2a_1=a_1q+a_1q^2\).
所以
\(q^2+q-2=0\).
解得
\(q=-2\)
（\(q=1\) 舍去）.
所以 \(\{a_n\}\) 的公比为 \(-2\).

（2）当 \(a_1=1\) 时，
\(a_n=\left( -2 \right)^{n-1}\)，
从而
\(na_n=n\left( -2 \right)^{n-1}\).
设数列 \(\{na_n\}\) 的前 \(n\) 项和为 \(S_n\)，则
\(S_n=1+2\left( -2 \right)+3\left( -2 \right)^2+\cdots+n\left( -2 \right)^{n-1}\)，
\(-2S_n=1\cdot \left( -2 \right)+2\cdot \left( -2 \right)^2+3\cdot \left( -2 \right)^3+\cdots+n\cdot \left( -2 \right)^n\).
两式相减得
\(3S_n=1+\left( -2 \right)+\left( -2 \right)^2+\cdots+\left( -2 \right)^{n-1}-n\left( -2 \right)^n\).
又
\(1+\left( -2 \right)+\left( -2 \right)^2+\cdots+\left( -2 \right)^{n-1}=\frac{1-\left( -2 \right)^n}{1-\left( -2 \right)}=\frac{1-\left( -2 \right)^n}{3}\)，
所以
\(3S_n=\frac{1-\left( -2 \right)^n}{3}-n\left( -2 \right)^n\).
故
\(S_n=\frac{1-\left( 1+3n \right)\left( -2 \right)^n}{9}\).
所以数列 \(\{na_n\}\) 的前 \(n\) 项和为
\(\frac{1-\left( 1+3n \right)\left( -2 \right)^n}{9}\).
\end{solution}

\begin{problem}
如图，\(D\) 为圆锥的顶点，\(O\) 是圆锥底面的圆心，\(AE\) 为底面直径，\(AE=AD\). \(\triangle ABC\) 是底面的内接正三角形，\(P\) 为 \(DO\) 上一点，\(PO=\dfrac{\sqrt{6}}{6}DO\).

\begin{enumerate}
\item 证明：\(PA\perp\) 平面 \(PBC\)；

\item 求二面角 \(B-PC-E\) 的余弦值.
\end{enumerate}

\bitmapfigure[max width=5.2cm]{img/2020/national_paper_1_science/q18_fig1.png}
\end{problem}

\begin{solution}
（1）不妨设圆 \(O\) 的半径为 \(1\)，则
\(OA=OB=OC=1\)，\(AE=AD=2\).
因为 \(\triangle ABC\) 是底面的内接正三角形，所以
\(AB=BC=AC=\sqrt{3}\).
又因为
\(DO=\sqrt{DA^2-OA^2}=\sqrt{3}\)，
所以
\(PO=\frac{\sqrt{6}}{6}DO=\frac{\sqrt{2}}{2}\).
从而
\(PA=PB=PC=\sqrt{PO^2+OA^2}=\frac{\sqrt{6}}{2}\).
在 \(\triangle PAC\) 中，
\(PA^2+PC^2=AC^2\)，
故
\(PA\perp PC\).
在 \(\triangle PAB\) 中，
\(PA^2+PB^2=2\times \left( \frac{\sqrt{6}}{2} \right)^2=3\).
又因为
\(AB=\sqrt{3}\)，
所以
\(PA^2+PB^2=AB^2\).
故
\(PA\perp PB\).
又因为 \(PB\cap PC=P\)，所以
\(PA\perp \text{平面 }PBC\).

（2）建立如图所示的空间直角坐标系，则 \(B\left( \frac{\sqrt{3}}{2},\frac{1}{2},0 \right)\)，\(C\left( -\frac{\sqrt{3}}{2},\frac{1}{2},0 \right)\)，\(P\left( 0,0,\frac{\sqrt{2}}{2} \right)\)，\(E\left( 0,1,0 \right)\).
所以 \(\overrightarrow{BC}=\left( -\sqrt{3},0,0 \right)\)，\(\overrightarrow{CE}=\left( \frac{\sqrt{3}}{2},\frac{1}{2},0 \right)\)，\(\overrightarrow{CP}=\left( \frac{\sqrt{3}}{2},-\frac{1}{2},\frac{\sqrt{2}}{2} \right)\).
设平面 \(PBC\) 的一个法向量为 \(\bs m=\left( x,y,z \right)\)，则 \(\bs m\cdot \overrightarrow{BC}=0\)，\(\bs m\cdot \overrightarrow{CP}=0\). 即 \(-\sqrt{3}x=0\)，\(\dfrac{\sqrt{3}}{2}x-\dfrac{1}{2}y+\dfrac{\sqrt{2}}{2}z=0\).
可取 \(\bs m=\left( 0,\sqrt{2},1 \right)\).
设平面 \(PCE\) 的一个法向量为 \(\bs n=\left( x,y,z \right)\)，则 \(\bs n\cdot \overrightarrow{CE}=0\)，\(\bs n\cdot \overrightarrow{CP}=0\). 即 \(\dfrac{\sqrt{3}}{2}x+\dfrac{1}{2}y=0\)，\(\dfrac{\sqrt{3}}{2}x-\dfrac{1}{2}y+\dfrac{\sqrt{2}}{2}z=0\).
两式相减可得 \(-y+\frac{\sqrt{2}}{2}z=0\)，即 \(z=\sqrt{2}y\).
由 \(\sqrt{3}x+y=0\) 得 \(y=-\sqrt{3}x\)，\(z=-\sqrt{6}x\).
取 \(x=\sqrt{2}\)，则 \(\bs n=\left( \sqrt{2},-\sqrt{6},-2\sqrt{3} \right)\).
故二面角 \(B-PC-E\) 的余弦值为 \(\dfrac{\left\lvert \bs m\cdot \bs n \right\rvert}{\left\lvert \bs m \right\rvert\left\lvert \bs n \right\rvert}=\dfrac{2\sqrt{5}}{5}\).
\end{solution}

\begin{problem}
甲，乙，丙三位同学进行羽毛球比赛，约定赛制如下：

累计负两场者被淘汰：比赛前抽签决定首先比赛的两人，另一人轮空；每场比赛的胜者与轮空者进行下一场比赛，负者下一场轮空，直至有一人被淘汰；当一人被淘汰后，剩余的两人继续比赛，直至其中一人被淘汰，另一人最终获胜，比赛结束.

经抽签，甲，乙首先比赛，丙轮空. 设每场比赛双方获胜的概率都为 \(\dfrac{1}{2}\).

\begin{enumerate}
\item 求甲连胜四场的概率；

\item 求需要进行第五场比赛的概率；

\item 求丙最终获胜的概率.
\end{enumerate}
\end{problem}

\begin{solution}
（1）甲连胜四场只能是前四场全胜，所以其概率为 \(\left( \frac{1}{2} \right)^4=\frac{1}{16}\).

（2）根据赛制，至少需要进行四场比赛，至多需要进行五场比赛.
比赛四场结束共有三种情况：甲连胜四场，乙连胜四场，丙上场后连胜三场. 其概率分别为 \(\frac{1}{16}\)，\(\frac{1}{16}\)，\(\frac{1}{8}\).
所以需要进行第五场比赛的概率为 \(1-\frac{1}{16}-\frac{1}{16}-\frac{1}{8}=\frac{3}{4}\).

（3）丙最终获胜有两种情况.
\(1^{\odot}\) 比赛四场结束且丙最终获胜，其概率为 \(\frac{1}{8}\).
\(2^{\odot}\) 比赛五场结束且丙最终获胜. 从第二场开始的四场比赛，按丙的胜，负，轮空结果有三种情况：\(\text{胜胜负胜}\)，\(\text{胜负空胜}\)，\(\text{负空胜胜}\)，其对应概率分别为 \(\frac{1}{16}\)，\(\frac{1}{8}\)，\(\frac{1}{8}\).
故丙最终获胜的概率为 \(\frac{1}{8}+\frac{1}{16}+\frac{1}{8}+\frac{1}{8}=\frac{7}{16}\).
\end{solution}

\begin{problem}
已知 \(A\)，\(B\) 分别为椭圆 \(E:\dfrac{x^2}{a^2}+y^2=1\)（\(a>1\)）的左，右顶点，\(G\) 为 \(E\) 的上顶点，\(\overrightarrow{AG}\cdot \overrightarrow{GB}=8\). \(P\) 为直线 \(x=6\) 上的动点，\(PA\) 与 \(E\) 的另一交点为 \(C\)，\(PB\) 与 \(E\) 的另一交点为 \(D\).

\begin{enumerate}
\item 求 \(E\) 的方程；

\item 证明：直线 \(CD\) 过定点.
\end{enumerate}
\end{problem}

\begin{solution}
如图示：
\bitmapfigure[max width=5.0cm]{img/2020/national_paper_1_science/q20_solution_fig1.png}
（1）由题意可得
\(A\left( -a,0 \right)\)，\(B\left( a,0 \right)\)，\(G\left( 0,1 \right)\).
所以
\(\overrightarrow{AG}=\left( a,1 \right)\)，\(\overrightarrow{GB}=\left( a,-1 \right)\).
由
\(\overrightarrow{AG}\cdot \overrightarrow{GB}=8\)
得
\(a^2-1=8\)，
故
\(a=3\).
所以椭圆 \(E\) 的方程为
\(\frac{x^2}{9}+y^2=1\).

（2）由（1）知
\(A\left( -3,0 \right)\)，\(B\left( 3,0 \right)\).
设
\(P\left( 6,m \right)\).
则直线 \(PA\) 的方程为
\(y=\frac{m}{9}\left( x+3 \right)\).
联立 \(\dfrac{x^2}{9}+y^2=1\) 和 \(y=\dfrac{m}{9}\left( x+3 \right)\)，得
\(\left( 9+m^2 \right)x^2+6m^2x+9m^2-81=0\).
由韦达定理可得
\(-3x_C=\frac{9m^2-81}{9+m^2}\)，
所以
\(x_C=\frac{-3m^2+27}{9+m^2}\).
代入直线 \(PA\) 的方程，可得
\(y_C=\frac{6m}{9+m^2}\).
即
\(C\left( \frac{-3m^2+27}{9+m^2},\frac{6m}{9+m^2} \right)\).
因为点 \(P\) 的坐标为 \(\left( 6,m \right)\)，点 \(B\) 的坐标为 \(\left( 3,0 \right)\)，所以直线 \(PB\) 的斜率为
\(\frac{m-0}{6-3}=\frac{m}{3}\).
故直线 \(PB\) 的方程为
\(y=\frac{m}{3}\left( x-3 \right)\).
联立 \(\dfrac{x^2}{9}+y^2=1\) 和 \(y=\dfrac{m}{3}\left( x-3 \right)\)，得
\(\left( 1+m^2 \right)x^2-6m^2x+9m^2-9=0\).
由韦达定理可得
\(3x_D=\frac{9m^2-9}{1+m^2}\)，
所以
\(x_D=\frac{3m^2-3}{1+m^2}\).
代入直线 \(PB\) 的方程，可得
\(y_D=-\frac{2m}{1+m^2}\).
即
\(D\left( \frac{3m^2-3}{1+m^2},-\frac{2m}{1+m^2} \right)\).
所以直线 \(CD\) 的斜率为
\(k_{CD}=\frac{y_C-y_D}{x_C-x_D}=\frac{4m}{3\left( 3-m^2 \right)}\).
于是直线 \(CD\) 的方程为
\(y+\frac{2m}{1+m^2}=\frac{4m}{3\left( 3-m^2 \right)}\left( x-\frac{3m^2-3}{1+m^2} \right)\).
化简得
\(y=\frac{4m}{3\left( 3-m^2 \right)}\left( x-\frac{3}{2} \right)\).
故直线 \(CD\) 过定点
\(\left( \frac{3}{2},0 \right)\).
\end{solution}

\begin{problem}
已知函数 \(f(x)=\e^x+ax^2-x\).

\begin{enumerate}
\item 当 \(a=1\) 时，讨论 \(f(x)\) 的单调性；

\item 当 \(x\ge 0\) 时，\(f(x)\ge \dfrac{1}{2}x^3+1\)，求 \(a\) 的取值范围.
\end{enumerate}
\end{problem}

\begin{solution}
（1）当 \(a=1\) 时，
\(f(x)=\e^x+x^2-x\).
所以
\(f'(x)=\e^x+2x-1\).
设
\(g(x)=f'(x)\)，
则
\(g'(x)=\e^x+2>0\).
所以 \(g(x)\) 在 \(\R\) 上递增，即 \(f'(x)\) 在 \(\R\) 上递增.
又因为
\(f'(0)=0\)，
所以当 \(x>0\) 时，\(f'(x)>0\)；当 \(x<0\) 时，\(f'(x)<0\).
故 \(f(x)\) 的增区间为 \(\left( 0,+\infty \right)\)，减区间为 \(\left( -\infty,0 \right)\).

（2）当 \(x=0\) 时，不等式恒成立.
当 \(x>0\) 时，由
\(\e^x+ax^2-x\ge \frac{1}{2}x^3+1\)
可得
\(a\ge \frac{\frac{1}{2}x^3+x+1-\e^x}{x^2}\).
设
\(h(x)=\frac{\frac{1}{2}x^3+x+1-\e^x}{x^2}\)，
则
\(h'(x)=\frac{\left( 2-x \right)\left( \e^x-\frac{1}{2}x^2-x-1 \right)}{x^3}\).
再设
\(m(x)=\e^x-\frac{1}{2}x^2-x-1\)，
则
\(m'(x)=\e^x-x-1\)，\(m''(x)=\e^x-1\).
因为 \(x\ge 0\) 时 \(m''(x)\ge 0\)，所以 \(m'(x)\) 在 \(\left( 0,+\infty \right)\) 上递增.
又因为
\(m'(0)=0\)，
所以 \(m'(x)\ge 0\)，故 \(m(x)\) 在 \(\left( 0,+\infty \right)\) 上递增.
又
\(m(0)=0\)，
所以
\(m(x)\ge 0\).
因此 \(h'(x)\) 的符号由 \(2-x\) 决定，即 \(h(x)\) 在 \(\left( 0,2 \right)\) 上递增，在 \(\left( 2,+\infty \right)\) 上递减.
所以
\(h_{\max}=h(2)=\frac{7-\e^2}{4}\).
故
\(a\ge \frac{7-\e^2}{4}\).
所以 \(a\) 的取值范围为
\(\left[ \frac{7-\e^2}{4},+\infty \right)\).
\end{solution}

\begin{problem}
在直角坐标系 \(xOy\) 中，曲线 \(C_1\) 的参数方程为 \(x=\cos^k t\)，\(y=\sin^k t\)（\(t\) 为参数）. 以坐标原点为极点，\(x\) 轴正半轴为极轴建立极坐标系，曲线 \(C_2\) 的极坐标方程为 \(4\rho\cos\theta-16\rho\sin\theta+3=0\).

\begin{enumerate}
\item 当 \(k=1\) 时，\(C_1\) 是什么曲线？

\item 当 \(k=4\) 时，求 \(C_1\) 与 \(C_2\) 的公共点的直角坐标.
\end{enumerate}
\end{problem}

\begin{solution}
（1）当 \(k=1\) 时，曲线 \(C_1\) 的参数方程为
\(x=\cos t\)，\(y=\sin t\)（\(t\text{ 为参数}\)）.
消去参数 \(t\)，可得
\(x^2+y^2=1\).
故 \(C_1\) 是以原点为圆心，以 \(1\) 为半径的圆.

（2）法一：当 \(k=4\) 时，曲线 \(C_1\) 的参数方程为
\(x=\cos^4 t\)，\(y=\sin^4 t\)（\(t\text{ 为参数}\)）.
消去参数 \(t\)，可得
\(\sqrt{x}+\sqrt{y}=1\).
由
\(4\rho\cos\theta-16\rho\sin\theta+3=0\)
并结合极坐标与直角坐标的互化公式，可得 \(C_2\) 的直角坐标方程为
\(4x-16y+3=0\).
联立 \(\sqrt{x}+\sqrt{y}=1\) 和 \(4x-16y+3=0\)，解得 \(x=\frac{1}{4}\)，\(y=\frac{1}{4}\).
故 \(C_1\) 与 \(C_2\) 的公共点的直角坐标为
\(\left( \frac{1}{4},\frac{1}{4} \right)\).

法二：当 \(k=4\) 时，曲线 \(C_1\) 的参数方程为
\(x=\cos^4 t\)，\(y=\sin^4 t\)（\(t\text{ 为参数}\)）.
两式作差可得
\(x-y=\cos^4 t-\sin^4 t=\cos^2 t-\sin^2 t=2\cos^2 t-1\)，
所以
\(\cos^2 t=\frac{x-y+1}{2}\)，
从而
\(x=\cos^4 t=\left( \frac{x-y+1}{2} \right)^2\).
整理得
\(\left( x-y \right)^2-2\left( x+y \right)+1=0\)（\(0\le x\le 1\)，\(0\le y\le 1\)）.
又因为
\(4\rho\cos\theta-16\rho\sin\theta+3=0\)，
且
\(x=\rho\cos\theta\)，\(y=\rho\sin\theta\)，
所以
\(4x-16y+3=0\).
联立 \(\left( x-y \right)^2-2\left( x+y \right)+1=0\) 和 \(4x-16y+3=0\)，解得 \(\begin{cases} x=\dfrac{169}{36},\\ y=\dfrac{49}{36} \end{cases}\)（舍去），或 \(\begin{cases} x=\dfrac{1}{4},\\ y=\dfrac{1}{4} \end{cases}\).
故 \(C_1\) 与 \(C_2\) 的公共点的直角坐标为
\(\left( \frac{1}{4},\frac{1}{4} \right)\).
\end{solution}

\begin{problem}
已知函数 \(f(x)=\left\lvert 3x+1 \right\rvert-2\left\lvert x-1 \right\rvert\).

\begin{enumerate}
\item 画出 \(y=f(x)\) 的图象；

\item 求不等式 \(f(x)>f(x+1)\) 的解集.
\end{enumerate}
\end{problem}

\begin{solution}
（1）函数
\examdisplaycases{f(x)=\left\lvert 3x+1 \right\rvert-2\left\lvert x-1 \right\rvert = }{
x+3, & x\ge 1,\\
5x-1, & -\dfrac{1}{3}\le x<1,\\
-x-3, & x<-\dfrac{1}{3}
}
图象如图所示：
\bitmapfigure[max width=5.3cm]{img/2020/national_paper_1_science/q23_solution_fig1.png}

（2）因为 \(f(x+1)\) 的图象是函数 \(f(x)\) 的图象向左平移 \(1\) 个单位所得，如图所示：
\bitmapfigure[max width=4.6cm]{img/2020/national_paper_1_science/q23_solution_fig2.png}
其中直线
\(y=5x-1\)
向左平移 \(1\) 个单位后可表示为
\(y=5\left( x+1 \right)-1=5x+4\).
联立 \(y=-x-3\) 和 \(y=5x+4\)，解得 \(x=-\frac{7}{6}\).
所以不等式 \(f(x)>f(x+1)\) 的解集为
\(\left\{ x\mid x<-\frac{7}{6} \right\}\).
\end{solution}
